\section{Auxiliary results independent of the Part II supposition}\label{app:partii}
The first paragraph of \src{Part II, p.85} supposes $\beta^*>7/8$ and defines its dynamic parameter $\kappa=2\beta^*-1$. Our uses of its auxiliary results require versions with this supposition removed. The following proposition specifies those versions and proves the removal by their mathematical dependencies.

\begin{proposition}\label{prop:partii}
The local statements in the following list remain valid with their original coefficient classes, zero extensions, bounded parameter ranges, masks, and uniformity, without the standing supposition $\beta^*>7/8$:
\[
\begin{gathered}
 \text{Lemmas 13.1--13.4; Corollary 14.1; Lemmas 14.2--14.3;}\\
 \text{Proposition 15.2; Proposition 16.1;}\\
 \text{Lemmas 17.1--17.3; Corollary 17.4; Lemmas 17.5--17.6;}\\
 \text{Lemma 18.1 in its printed range, Lemmas 18.2--18.3; Lemma 19.1.}
\end{gathered}
\]
For positive slots, Lemma~18.1 retains its explicit premise $\beta^*\le(1+\kappa)/2$ when $\kappa<1$. Proposition~16.1 retains the explicit primitive conjugate-numerator bound \eqref{eq:correctnumerator}. Lemma~19.1 retains its buffered and cumulative-frequency hypotheses. All other results in this list need no assertion about the location of Hecke zeros.
\end{proposition}
\begin{proof}
\emph{Finite correlations and marked reflection.}
Lemma~13.1 is the prime ideal asymptotic in a fixed ray quotient, with a fixed annular test and fixed finite deletions. Its proof uses the prime ideal theorem in a fixed Frobenius class \src{\S13.2} and Stieltjes integration by parts; neither the quotient nor the deletion set moves with the prime scale. Its proof does not use the family supremum.
Lemma~13.2 evaluates prime-power Fourier sums. Lemma~13.3 is finite Fourier inversion of two complete Gauss sums, retaining their common factors and the original principal zero convention. Lemma~13.4 is its CRT factorization after the complete shared prime support is extracted; the residual factors are genuinely coprime there. These are finite identities. They apply to the whole products and exact masks retained in Sections~\ref{sec:low} and \ref{sec:fourth}.

Corollary~14.1 applies the Part~I completed cubic reflection, \src{Proposition 5.1}, to the product-form whole-index mark. Its conditions are a fixed finite ray family, disjoint fixed prime lists, coefficients independent of the row and fourth-power label, and the original puncture. Expand each local $1-\chi_\p^0$. The inactive coefficient difference is $q_\p^{-1}$, an active mark has exponent zero, and the row exponent is one. CRT gives the row-slot cubic exponent $1+0+2=3$, whose phase is one. The remaining two zero-preserving symbols are those of \src{(14.16)}. This reasoning contains no reciprocal of a Hecke function and no geometry.

Lemma~14.2 uses the squarefree indices $k,n,P$, positive norm scales $K,N,B,L\ge1$, and two coefficient-independence conditions: $a(P)$ is independent of $k,n,b$ and $\beta(n,b)$ is independent of $k,P$. The indices $k,P$ avoid $S$; the source $n,b$ may contain primes of $S$ other than $\lambda$. Expand the full $(P,k)$ mask, write $P=DP'$, $k=Dk'$, and freeze the other divisor and row labels before applying the quadratic reduction \src{(5.27)} and the cubic norm \src{(14.6)}. In the cubic step the squarefree column $j=c'm=n/r$ may contain primes of $S$ other than $\lambda$; \src{p.96} passes it to \src{(14.6)} after cubic reciprocity and conjugation, since (14.6) excludes no primes. Its output is
\[
 (K+NB)B\{N+L+(NL)^{2/3}\},
\]
up to small powers. The residual divisor factors only improve this expression. There is no zero-location premise.

Lemma~14.3 combines these two facts with the common kernel and lattice tails of \src{Lemmas 5.3--5.4}. It is applied on actual residual-row dyads, with bounded lengths, fixed characters and punctures, disjoint product-form slots, and row restrictions independent of active slots and dual variables except for their explicit coprimality mask. The exact joint kernel is \src{(14.15)}, including every dual and slot norm. Its squared scalar saving is $(T_d-v-3\ell_b-e_\lambda)_+/2$. The hybrid norm gives \eqref{eq:Eref}; summing the fixed powerful and supported row parts adds $O/2$ once. The kernel contours are those of the completed theta reflection, not reciprocal Hecke contours. Thus neither the standing supposition nor the geometry \src{(12.4)} enters this proof.

\emph{The additive norm and local errors.}
Proposition~15.2 assumes fixed $S,T,b_{\rm aux},\xi,W_1,c_T$, polynomial $Q,Y'\ge1$, $P_a=Y'^2/Q\ge1$, a fixed row ball, and precisely its finite ray-coefficient class. Expanding its square and applying Lemma~13.3 gives the full common-factor correlation. Freeze the fixed ray parts before periodicizing the joint coefficient. On nonexceptional frequencies a nontrivial good-prime character has zero complete mean. Four-dimensional lattice Poisson gives the bound in \src{(15.5)}, and summing the exact divisor and row counts yields $(Q/Y')(P_a^2/Y')$. Exceptional frequencies have $(k)=a_1^6e$, with $e$ supported on the extracted factor and fixed set; elementary counting gives $(Q/Y')P_a^{1/6}$. The zero frequency is the original diagonal and contributes $Q/Y'$. These are exactly the three terms in \src{(15.4)}. No use is made of the geometry-dependent Proposition~15.3.

Proposition~16.1 assumes its displayed real strip, sixth-power-free physical rows outside $S$, nonprincipal primitive numerator, nonzero central selected local factors, and \eqref{eq:correctnumerator}. Its proof uses the six-case table \eqref{eq:localtable} and one conductor allocation to the entire set of strict ramified primes. The functional-equation conductor exponent is $a-1/2+6e$; its validity requires $51/100\le a\le1$, not $a>7/8$. Section~\ref{sec:geometry} establishes its explicit conjugate premise by the common buffer. Lemma~\ref{lem:cutoff} supplies its cutoff without a target-dependent choice. The separate principal-factor paragraph \src{(16.12)--(16.14)} is excluded; Proposition~\ref{prop:correction} proves its replacement on $\mathcal D'_2$.

\emph{Finite propagation and masked Poisson.}
Lemma~17.3 and Corollary~17.4 are finite backward propagation on an acyclic graph. Their hypotheses are fixed real terminal and edge exponents, finite weighted moments of one common coefficient measure, fixed-dimensional smooth profile bounds, and fixed linear height maps. Normalizing by the prescribed norm exponent makes every label measure have bounded small-power mass. Backward induction raises finite seminorm and height orders but leaves the real norm powers fixed. All these statements concern nonnegative integrals and finite graph depth; no zero-location assertion occurs.

Lemma~17.5 is masked primitive Poisson. With a primitive finite character $\psi$ modulo $\mathfrak m$, an ideal $R$, a radial Schwartz test, and $K>0$, its formula is
\[
 \sum_k\psi(k)\1_{(k,R)=1}\Phi(q_k/K)
 =\frac{K\gamma(\psi;\mathfrak m)}{\sqrt{q_{\mathfrak m}}}
   \sum_{d\mid\rad R}\frac{\mu(d)\psi(d)}{q_d}
   \sum_h\overline{\psi(h)}
          \mathcal F\Phi\!\left(\frac{Kq_h}{q_dq_{\mathfrak m}}\right).
\]
Expand the complete mask, substitute $k=dk'$, and use finite Fourier inversion modulo $\mathfrak m$. The finite transform is $\sqrt{q_{\mathfrak m}}\gamma(\psi;\mathfrak m)\overline{\psi(h)}$. Primitivity kills nonunit frequencies; Parseval gives $|\gamma|=1$ for a nonprincipal primitive character. For the primitive principal modulus use $\mathfrak m=1$, $\gamma=1$, and value one at zero. At each use here the row scale is at least one. Its principal nonzero-frequency restoration costs $O(KZ^\eps)$ by lattice counting and Schwartz decay, with the complete outer mask retained. This is a finite-arithmetic and lattice statement.

\emph{Canonical marked estimate and its initialization.}
For precision the canonical family of Lemma~17.2 has fixed puncture $\rho(n)=\1_{(n,r_\rho)=1}$, fixed finite-ray $\nu$, disjoint product-form prime mark $d(n)$ of total length at most $z_0$, and squarefree labels $q_f\asymp Z^V$ with a nonnegative divisor-bounded weight $w(f)$ depending only on $f$. Put $E=N+V$. Its two local premises are
\[
 q_{r_\rho}\le Z^{E-M-z_0-c_*},\qquad
 3M+6z_0+c_*\le4E,\qquad c_*>0.
\]
The normalized family is
\[
 Z^{-V}\sum_f w(f)\sum_{0<q_k\ll Z^M}
 \left|Z^{-N/2}\sum_{n\ \sfv}
       \overline{\alpha(n)}\gamma_2(n)\nu(n)\rho(n)
       \chi_n(k)\chi_n(f)^4d(n)W(q_n/Z^N)\right|^2.
\]
No additional arbitrary column coefficient is allowed.

Its terminal proof uses the exact cube-Möbius completion \src{(17.10)} and the assumption-free Lemma~14.3. The node invariants are
\[
 E-M-Q-z_0\ge c_{\rm node},\qquad4E-3M-6z_0\ge c_{\rm node},
 \quad Q=\log_Zq_{r_\rho}.
\]
For step $d\le c_*/200$, $c_{\rm node}\ge c_*/2$, and support and local losses at most $c_*/1000$, the three reflected branches \src{(17.13)--(17.14)} are strictly below $E$:
\[
 E-c_{\rm node}+5\eta/3,\quad
 E-2c_{\rm node}+5d+17\eta/3+\tau_{\rm ref},\quad
 E-M/4-3c_{\rm node}/4+5d/2+19\eta/6+\tau_{\rm ref}/2.
\]
The complete weight in the $f$-sum cancels its outside $Z^{-V}$; no extra label count is inserted into the hybrid norm.

The remaining canonical terms are reduced by the two finite masked Poisson transforms \src{(17.15)--(17.65)}. Priority slot allocation is an exact partition even at shared primes. The first transform is used on genuine coprime columns, its raw tail is removed before full mutual-gcd Möbius extension, and its kernel is one joint profile with all norm roots retained. The old $f$ label is removed once by the divisor-fiber count \src{(17.32)}. The second transform reconstructs a canonical squarefree label $f_{\rm new}=JCd_2v'$; its squarefreeness follows from the retained coefficient zeros, not a new mask. Its puncture is fixed within that child, independent of its varying row and label. The child margins in \src{(17.51)} lose at most $5\eta$ and $7\eta$, and its row width drops at least $2d-8\eta$. The exact normalization \src{(17.64)} is
\[
 \kappa_i+\lambda_c+C_c+2F_c=E,
\]
where $\kappa_i$ is a normalization exponent, not a moment parameter. With $\eta\le d/16$, the per-edge loss is $40\eta+\tau+\pi$. Choose the finite depth and then $\eta,\tau,\pi$ so that the sum of these losses is below the requested $\eps$, as in \src{(17.66)}. All raw tail orders are chosen after the crude count, and the finite propagation rule fixes internal orders before an external cutoff. These invariants and transformations contain neither $\beta^*$ nor a dynamic-bin inequality.

Lemma~17.1 follows by its single masked-Poisson initialization \src{(17.67)--(17.86)}. The exact overlap split is $n=jn_0$, $P=jP_0$; the scalar $Z^{-G}$ pays the assigned prime count. The resulting canonical label is squarefree and its full puncture remains within its local premise. For prescribed inverse margins $c_1,c_2>0$, the child margins are at least $c_1-9\eta-\tau_{\rm init}$ and $c_2-22\eta-3\tau_{\rm init}$, so $c_*=\min(c_1,c_2)/2$ is allowed after small losses. The energy on each side is $m+11\eta+\pi_{\rm init}+\eps_c$. Cauchy averages rather than doubles these costs. This derives the marked inverse estimate under its stated strict margins with no family-boundary input.

Finally \src{(17.87)--(17.90)} use its zero-slot specialization, scale Sobolev, and the exact zero-preserving divisor identity. On sixth-power-free rows the map $(u,a)\mapsto ua^6$ is injective. This proves Lemma~17.6 under its printed $H=\max(2U,D^{1+c})$ and bounded profile ranges, again without a zero-location hypothesis.

\emph{The fourth moment and prime amplitudes.}
For Lemma~18.1, simultaneous mask deletion and primitive reflection \src{(18.4)--(18.10)} use only the general smooth calculus, functional equation, and deleted-factor estimates. Its sole zero-location input is \src{(18.11)} on $s_\kappa=(1+\kappa)/2\ge\beta^*$. It is absent for zero slots and replaced by absolute prime counting at $\kappa=1$. The terminal widths, same-band comparison, two complete finite transforms, paired clipping, and exceptional cancellation \src{(18.13)--(18.50)} use only the local length inequalities and $\kappa$. The finite-depth choices \src{(18.51)--(18.52)} close the estimates. Section~\ref{sec:fourth} rederives every inequality affected by extending that parameter below $3/4$.

Lemma~18.2 is preservation of the complete common coefficient under those transforms and full Möbius allocations; Lemma~18.3 is lattice Poisson with one fixed character, one common polynomial-norm squarefree mask, one norm power per side, and the four original formal scales. Their equal-product leading terms cancel, and their absolute-volume fallback is valid at every positive scale. These results have no zero-location premise.

Lemma~19.1 is a buffered prime-annulus consequence of Part~I Lemma~8.2 and Lemma~4.9. The prime-annulus frequency is within the same cumulative allowance, and the main slot retains its original identity-ray restriction and zero extension. Its logarithmic derivative is controlled only on that buffer; its exterior tails stay on the absolute line. Thus its only analytic premise is the stated buffer, not the standing contradiction.

This establishes each item of the list with exactly the retained local hypotheses.
\end{proof}

The fixed-geometry conclusion of \src{Proposition 15.3}, the positive-gap capacity replacement in \src{Proposition 19.2}, and \src{Proposition 20.3} are not inputs here. Nor are the fixed-boundary conclusions of \src{Lemmas 10.3--10.6} or \src{Lemma 16.2} used outside their printed ranges. Their needed parameter-dependent estimates were derived in Sections~\ref{sec:low}--\ref{sec:cubic}. The pure identity in \src{Lemma 20.2} was reproduced in \eqref{eq:oldpoly}--\eqref{eq:oldsquare}.

The paragraph \emph{The principal local factor} of \src{\S16.2}, namely \src{(16.12)--(16.14)}, is not an input either: it is stated on the second Euler region \src{(7.15)} ($\Rea s\ge7/8$) and is replaced by Proposition~\ref{prop:correction} on $\mathcal D'_2$.
